% 第12章 课程综合项目与创新实践
% 智慧银行实验教程——AI驱动的金融科技实践

\chapter{第12章\quad 课程综合项目与创新实践}
\chapteroverview{
本章学习目标：
\begin{enumerate}
  \item 掌握综合项目的选题方法
  \item 学会项目管理与团队协作
  \item 了解成果展示与答辩规范
  \item 能将实验成果转化为学术论文
  \item 了解金融科技竞赛参赛路径
\end{enumerate}
}

\keyterms{项目选题、可行性分析、WBS工作分解、Git协作、答辩演示、学术成果转化、金融科技竞赛、挑战杯、互联网+}

% ============================================================
\section{项目选题指南}
% ============================================================

\subsection{选题原则}

综合项目是本课程的"收官之作"，选题质量直接决定项目的成败。好的选题应满足四个原则：

\textbf{可行性（Feasibility）}：项目必须在4周内、2--4人团队、使用课程所学工具可完成。避免选择需要海量数据、复杂硬件或超长开发周期的项目。判断标准：能否在1周内搭建出最小可用原型（MVP）。

\textbf{创新性（Novelty）}：项目应在某个维度上有所创新——可以是技术上的新组合（如LLM+知识图谱用于银行问答）、业务上的新场景（如ESG评级工具）、或方法上的新尝试（如BMAD方法论用于金融系统设计）。但创新不等于"从零发明"，"已有技术的跨界应用"同样是创新。

\textbf{金融相关性（Finance Relevance）}：项目必须与金融/银行领域紧密相关。纯技术项目（如"通用聊天机器人"）不符合课程定位。好的选题应能回答"这个项目解决了什么金融问题"。

\textbf{技术适配性（Tech Fit）}：项目应能充分发挥AI IDE + MCP + Skill的技术栈优势。如果项目完全不需要AI参与，说明没有用上课程核心能力。

\begin{notebox}[选题自查清单]
在确定选题前，请逐项检查：
\begin{itemize}[leftmargin=1em]
  \item [\checkmark] 4周内能否完成MVP？
  \item [\checkmark] 是否有至少一个创新点？
  \item [\checkmark] 解决了什么金融问题？
  \item [\checkmark] 是否需要AI参与？如何参与？
  \item [\checkmark] 数据从哪里来？
  \item [\checkmark] 技术难点是什么？能否克服？
\end{itemize}
\end{notebox}

\subsection{可选项目方向}

以下提供12个可选项目方向，每个方向附带难度评级和简要介绍：

\textbf{1. 银行智能网点导航系统（★★）}

基于银行网点信息，构建一个智能导航系统，用户输入位置和需求（如"办理信用卡"），系统推荐最近的网点并显示等候时间估算。技术栈：HTML+JS+地图API，可使用Playwright MCP抓取网点信息。适合2人团队，2周可完成MVP。

\textbf{2. 个人财务健康诊断助手（★★）}

用户输入月收入、支出、负债等数据，系统从储蓄率、负债率、流动性等维度评估财务健康状况，生成诊断报告和改善建议。技术栈：HTML+JS+图表库（ECharts），AI辅助生成评估逻辑和理财建议。适合2人团队，2周可完成。

\textbf{3. 银行反欺诈交易检测系统（★★★）}

基于模拟的信用卡交易数据，使用规则引擎+机器学习模型检测异常交易（大额异地消费、频繁小额试探等）。技术栈：Python+scikit-learn+Flask前端，需构造模拟数据集。适合3人团队，3周可完成。

\textbf{4. 智能贷款计算器与方案比较（★★）}

支持等额本息、等额本金、先息后本等多种还款方式计算，可比较不同银行、不同期限的贷款方案，生成可视化对比报告。技术栈：HTML+JS+ECharts，AI辅助编写计算逻辑和交互设计。适合2人团队，2周可完成。

\textbf{5. 银行客户流失预警系统（★★★）}

基于模拟的客户行为数据（交易频率下降、余额持续走低、投诉增加等），构建客户流失预警模型，识别高流失风险客户并推荐挽留策略。技术栈：Python+scikit-learn+Flask，需构造特征工程。适合3人团队，3周可完成。

\textbf{6. 数字人民币模拟交易系统（★★★）}

模拟数字人民币（e-CNY）的开通、充值、转账、消费全流程，包含双离线支付的基本逻辑演示。技术栈：HTML+JS+localStorage，重点在于业务流程的正确模拟。适合3人团队，3周可完成。

\textbf{7. 银行ESG评级分析工具（★★★）}

基于公开ESG数据，为银行客户提供ESG评级查询和对比分析功能。包含环境（E）、社会（S）、治理（G）三个维度的评分展示。技术栈：Python+Flask+ECharts，数据来源为模拟或公开ESG报告。适合3人团队，3周可完成。

\textbf{8. 跨境汇款路径优化系统（★★★★）}

根据汇率、手续费、到账时间等维度，为用户推荐最优的跨境汇款方案。需对接多个汇款渠道的数据（可模拟），并考虑外汇管制因素。技术栈：Python+Flask+汇率API，逻辑复杂度高。适合4人团队，4周可完成。

\textbf{9. 银行知识图谱构建与问答（★★★★）}

基于银行产品、业务流程、监管法规等知识构建知识图谱，实现基于图谱的智能问答。技术栈：Python+Neo4j（或JSON模拟）+LLM，需设计本体和实体关系。适合4人团队，4周可完成。

\textbf{10. 中小银行数字化转型评估工具（★★★）}

参考银保监会的数字化转型监管要求，设计一套评估指标体系，用户输入银行各项能力数据后自动生成转型成熟度评估报告。技术栈：HTML+JS+ECharts，重点在于指标体系设计。适合3人团队，3周可完成。

\textbf{11. 供应链金融平台原型（★★★★）}

模拟核心企业—供应商—银行三方参与的供应链金融平台，包含应收账款融资、订单融资等场景。需实现多角色登录和业务流程串联。技术栈：HTML+JS+localStorage+Flask，业务逻辑复杂。适合4人团队，4周可完成。

\textbf{12. 银行合规检查自动化系统（★★★）}

将银行常见合规检查项（反洗钱、投资者适当性、信息披露等）结构化为检查清单，AI辅助自动审查文档合规性，生成合规报告。技术栈：HTML+JS+AI对话，需设计检查规则库。适合3人团队，3周可完成。

\begin{notebox}[选题建议]
\begin{itemize}[leftmargin=1em]
  \item 初学者推荐选择★★项目，2人团队即可完成
  \item 有编程基础的同学可选择★★★项目，更具挑战性
  \item ★★★★项目适合有全栈开发经验的4人团队
  \item 也欢迎自行设计选题，但需经教师审批
  \item 无论选择哪个方向，务必在第一周完成MVP
\end{itemize}
\end{notebox}

% ============================================================
\section{项目管理与团队协作}
% ============================================================

\subsection{团队角色分工}

2--4人团队建议按以下角色分工（可兼任）：

\begin{table}[H]
\centering
\begin{tabular}{llp{6cm}}
\toprule
\textbf{角色} & \textbf{人数} & \textbf{职责} \\
\midrule
项目经理 & 1人 & 需求分析、任务分解、进度跟踪、最终答辩主讲 \\
前端开发 & 1--2人 & 界面设计、交互实现、用户体验优化 \\
后端/数据 & 1人 & 数据处理、算法实现、API开发 \\
文档/测试 & 1人（可兼任） & 需求文档、测试用例、实验报告撰写 \\
\bottomrule
\end{tabular}
\caption{团队角色分工}
\end{table}

2人团队建议：一人兼项目经理+前端，一人兼后端+文档。3人团队建议：项目经理兼前端，一人后端，一人文档测试。4人团队可完整覆盖四个角色。

\subsection{4周项目时间线}

\begin{table}[H]
\centering
\small
\begin{tabular}{clp{7cm}}
\toprule
\textbf{周次} & \textbf{阶段} & \textbf{关键任务} \\
\midrule
第1周 & 规划 & 选题确认、需求分析、MVP定义、技术选型、Git仓库初始化、AI协作计划 \\
第2周 & 开发（核心） & MVP开发、核心功能实现、每日站会（5分钟）、AI辅助编码 \\
第3周 & 开发（完善） & 功能完善、边界处理、UI美化、性能优化、Bug修复、内部测试 \\
第4周 & 答辩 & 演示视频录制、答辩PPT制作、实验报告撰写、代码整理、最终提交 \\
\bottomrule
\end{tabular}
\caption{4周项目时间线}
\end{table}

\textbf{第1周详细计划}：

\begin{enumerate}[leftmargin=2em]
  \item Day 1--2：选题讨论与确认，完成项目立项书（1页）
  \item Day 3：需求分析与功能列表，用AI生成WBS
  \item Day 4：技术选型与架构设计，确定技术栈
  \item Day 5：Git仓库初始化，创建项目脚手架
  \item Day 6--7：MVP核心功能开发
\end{enumerate}

\textbf{关键里程碑}：

\begin{itemize}[leftmargin=2em]
  \item Day 3：选题锁定，不得再更换
  \item Day 7：MVP可用，能演示核心流程
  \item Day 14：功能完整，基本无Bug
  \item Day 21：全部功能就绪，开始准备答辩
  \item Day 28：最终提交
\end{itemize}

\subsection{Git协作工作流}

团队项目必须使用Git进行版本控制和协作。推荐采用Feature Branch工作流：

\begin{lstlisting}[style=shell]
# 初始化仓库
git init
git remote add origin https://cnb.cool/team/project.git

# 创建主分支
git checkout -b main

# 每个功能在独立分支开发
git checkout -b feature/login-page
# ... 开发 ...
git add .
git commit -m "feat: 完成登录页面"
git push origin feature/login-page

# 合并到main（通过Pull Request）
# 在CNB平台上创建PR，AI辅助代码审查后合并
\end{lstlisting}

\textbf{分支命名规范}：

\begin{table}[H]
\centering
\begin{tabular}{ll}
\toprule
\textbf{分支类型} & \textbf{命名示例} \\
\midrule
功能分支 & feature/login-page, feature/faq-search \\
修复分支 & fix/chat-scroll-bug \\
文档分支 & docs/api-reference \\
发布分支 & release/v1.0 \\
\bottomrule
\end{tabular}
\caption{Git分支命名规范}
\end{table}

\textbf{提交信息规范}（Conventional Commits）：

\begin{lstlisting}
feat: 新增FAQ知识库检索功能
fix: 修复聊天窗口滚动不到底的问题
docs: 更新API接口文档
style: 调整按钮间距和颜色
refactor: 重构对话流引擎
test: 添加检索引擎单元测试
\end{lstlisting}

\subsection{AI在项目管理中的应用}

AI不仅能辅助编码，还能在项目管理全流程中发挥作用：

\textbf{任务分解}：用AI生成WBS（工作分解结构）。

\begin{practicebox}[用AI生成项目WBS]
\textbf{提示词}：

\begin{lstlisting}
我要开发一个[项目名称]，团队2--4人，4周时间。
请帮我生成WBS工作分解结构，要求：
1. 按周分解（4周），每周列出具体任务
2. 每个任务标注：任务名、负责人角色、预计工时、依赖关系
3. 标注关键路径上的任务
4. 输出为Markdown表格格式
\end{lstlisting}
\end{practicebox}

\textbf{进度跟踪}：用AI自动生成日报/周报。

\begin{lstlisting}
请根据以下Git提交记录，生成本周项目进展报告：
- feat: 完成FAQ知识库50条数据录入
- feat: 实现关键词检索引擎
- fix: 修复消息气泡显示错位问题
- docs: 编写API接口文档

报告格式：本周完成、下周计划、风险与阻塞
\end{lstlisting}

\textbf{代码审查}：AI辅助Review Pull Request。

在CNB平台上创建PR后，可使用AI辅助审查代码质量、安全问题和规范一致性。审查要点包括：变量命名是否规范、是否有硬编码的敏感信息、错误处理是否完善、代码逻辑是否清晰。

% ============================================================
\section{成果展示与答辩规范}
% ============================================================

\subsection{演示文稿结构}

答辩演示应遵循"问题→方案→演示→技术→总结"五段式结构：

\begin{table}[H]
\centering
\begin{tabular}{clp{6cm}}
\toprule
\textbf{环节} & \textbf{时间} & \textbf{内容} \\
\midrule
问题 & 30秒 & 我们要解决什么金融问题？为什么重要？ \\
方案 & 1分钟 & 我们的解决方案是什么？创新点在哪？ \\
演示 & 2分钟 & 现场演示系统核心功能（最震撼的部分） \\
技术 & 1分钟 & 技术架构、AI协作方式、关键技术决策 \\
总结 & 30秒 & 成果亮点、不足与改进方向、团队分工 \\
\bottomrule
\end{tabular}
\caption{5分钟答辩时间分配}
\end{table}

\subsection{Demo演示技巧}

Demo演示是答辩的核心环节，以下是关键技巧：

\textbf{提前录制备用视频}。现场Demo可能因网络、设备等原因失败，务必提前录制完整的演示视频作为备选方案。录制时确保画面清晰、操作流畅、讲解同步。

\textbf{从用户视角出发}。不要先展示技术架构，而是先展示用户能体验到的功能。"让评委先看到效果，再了解技术"——这比"先讲原理再展示"更吸引人。

\textbf{准备多个演示场景}。主场景用于正式演示，备选场景用于回答评委提问时的即时展示。例如主场景演示正常流程，备选场景演示异常处理。

\textbf{预填充数据}。确保演示系统中有足够的预填充数据，避免"空壳"展示。例如贷款计算器应预设多种贷款方案，客服系统应有完整的FAQ和对话记录。

\subsection{答辩常见问题准备}

根据往届答辩经验，评委常问的问题包括：

\begin{enumerate}[leftmargin=2em]
  \item 你们的创新点是什么？与市面上已有的产品有什么区别？
  \item AI在你的项目中扮演什么角色？如果没有AI会怎样？
  \item 你们的系统如何保证数据安全和合规性？
  \item 如果给你更多时间，你会做哪些改进？
  \item 项目的最大技术难点是什么？如何解决的？
  \item 用户体验上有什么考量？如何收集用户反馈？
  \item 你们的模型准确率是多少？如何评估？
\end{enumerate}

\begin{notebox}[答辩技巧]
\begin{itemize}[leftmargin=1em]
  \item 不会的问题诚实回答"这是一个很好的问题，我们目前还没有深入考虑，后续会重点关注"
  \item 不要与评委争论，可以委婉表达不同观点
  \item 团队成员都应参与回答，展示团队协作
  \item 控制时间，超时比内容少更影响评分
\end{itemize}
\end{notebox}

% ============================================================
\section{从实验到论文：学术成果转化}
% ============================================================

\subsection{实验报告与学术论文的区别}

课程实验报告和学术论文在目标、读者和结构上有显著差异：

\begin{table}[H]
\centering
\begin{tabular}{lll}
\toprule
\textbf{维度} & \textbf{实验报告} & \textbf{学术论文} \\
\midrule
目标 & 证明完成了实验 & 贡献新知识或新方法 \\
读者 & 任课教师 & 学术同行 \\
结构 & 按实验步骤组织 & 引言→文献→方法→实验→结论 \\
创新要求 & 不要求 & 必须有创新点 \\
文献引用 & 可选 & 必须充分引用 \\
数据要求 & 可用模拟数据 & 需要真实数据或严谨实验 \\
篇幅 & 3000--5000字 & 6000--15000字 \\
\bottomrule
\end{tabular}
\caption{实验报告与学术论文对比}
\end{table}

将实验成果转化为学术论文，关键在于找到"新"的角度——新技术在老场景中的应用、新方法对老问题的解决、新工具对教学效果的改善等。

\subsection{论文选题方向}

结合本课程内容，推荐以下三类论文选题方向：

\textbf{技术应用类}：聚焦AI技术在金融场景中的具体应用。典型选题：

\begin{itemize}[leftmargin=2em]
  \item 「基于LLM Agent的银行CRM系统设计与实现」
  \item 「MCP协议在银行智能客服系统中的应用研究」
  \item 「基于Skill体系的金融数据分析工作流构建」
  \item 「大语言模型驱动的银行FAQ知识库自动构建方法」
  \item 「AI辅助的银行合规检查自动化方案设计」
\end{itemize}

\textbf{实证研究类}：通过实验数据验证AI工具的效果。典型选题：

\begin{itemize}[leftmargin=2em]
  \item 「AI辅助编程工具对金融专业学生编程能力的影响——基于XX实验的实证研究」
  \item 「大语言模型在银行FAQ检索中的准确率评估」
  \item 「BMAD方法论对金融信息系统开发效率的影响研究」
  \item 「AI IDE对不同编程水平学生的差异化影响分析」
  \item 「对话式AI与图形界面在银行业务办理中的用户体验比较」
\end{itemize}

\textbf{方法论类}：提出新的方法框架或改进方案。典型选题：

\begin{itemize}[leftmargin=2em]
  \item 「BMAD方法论在金融信息系统开发中的应用与改进」
  \item 「银行智能客服的多轮对话流设计方法研究」
  \item 「基于MCP协议的金融工具集成框架设计」
  \item 「面向金融专业学生的AI协作能力培养路径研究」
  \item 「银行合规话术的自动生成与过滤方法」
\end{itemize}

\subsection{论文结构指导}

一篇完整的学术论文应包含以下结构：

\textbf{引言（Introduction）}：阐述研究背景、问题描述、研究意义和文章结构。引言要回答三个问题：为什么要研究这个？别人做到了什么程度？我要做什么？

\textbf{文献综述（Literature Review）}：系统梳理相关领域的研究现状，指出现有研究的不足，从而引出本文的创新点。文献引用不少于15篇，其中近3年文献不少于50\%。

\textbf{方法（Methodology）}：详细描述所采用的方法、技术架构和实现方案。可包含系统架构图、算法伪代码、对话流设计图等。方法是论文的"贡献"部分，必须写得足够详细，使他人可以复现。

\textbf{实验（Experiments）}：描述实验设置、数据集、评估指标和实验结果。实验应包含对照组或基线方法，以证明方法的有效性。结果应以表格或图表形式清晰展示。

\textbf{讨论（Discussion）}：分析实验结果，讨论方法的优势和局限性，与已有方法进行比较，提出改进方向。

\textbf{结论（Conclusion）}：总结研究贡献，明确创新点，指出局限性和未来工作方向。

\subsection{投稿目标}

根据论文类型和质量，推荐以下投稿目标：

\begin{table}[H]
\centering
\small
\begin{tabular}{llp{5cm}}
\toprule
\textbf{论文类型} & \textbf{推荐目标} & \textbf{说明} \\
\midrule
技术应用类 & 《计算机工程与应用》《计算机科学》 & 中文核心期刊，接收系统设计类论文 \\
实证研究类 & 《实验技术与管理》《实验室研究与探索》 & 教育类核心期刊，接收教学实验研究 \\
方法论类 & 《金融科技时代》《金融电子化》 & 金融科技行业期刊，接收方法论创新 \\
高质量论文 & 《软件学报》《计算机学报》 & CCF推荐期刊，要求高但影响力大 \\
会议论文 & 各类金融科技/人工智能学术会议 & 审稿周期短，适合快速发表 \\
\bottomrule
\end{tabular}
\caption{投稿目标参考}
\end{table}

% ============================================================
\section{金融科技竞赛指南}
% ============================================================

\subsection{国内主要竞赛}

金融科技领域有以下主要竞赛可供参与：

\textbf{1. "挑战杯"大学生课外学术科技作品竞赛}

"挑战杯"是全国最具影响力的大学生学术科技竞赛，每两年举办一次（奇数年国赛）。竞赛分为自然科学类学术论文、哲学社会科学类调查报告、科技发明制作三大类别。金融科技项目通常归入科技发明制作类或哲学社会科学类。

\textbf{参赛要求}：作品需有完整的学术报告和实物/软件展示，团队不超过8人，需有指导教师。

\textbf{2. 全国大学生金融科技创新大赛}

由中国金融教育发展基金会等主办，聚焦金融科技领域的创新应用。竞赛设置银行数字化转型、智能风控、普惠金融、绿色金融等赛道。作品形式为可运行的软件系统+完整的商业计划书。

\textbf{参赛要求}：团队3--5人，需提交作品演示视频和技术文档。

\textbf{3. "互联网+"大学生创新创业大赛（金融赛道）}

由教育部主办，是目前规模最大的大学生创新创业竞赛。竞赛设置高教主赛道（创意组/初创组/成长组）和产业命题赛道。金融科技类项目通常参加高教主赛道或产业命题赛道的金融命题。

\textbf{参赛要求}：团队3--15人，需提交商业计划书和路演视频。

\textbf{4. 各省"金融科技"创新应用大赛}

各省金融监管局或金融学会主办，面向金融机构和高校团队。竞赛内容通常围绕金融科技在普惠金融、风险防控、监管科技等领域的创新应用。部分省份设有高校专项赛道。

\textbf{5. 中国大学生计算机设计大赛}

由教育部高校计算机类专业教指委主办，设有大数据、人工智能、信息可视化等类别。金融数据分析、智能客服等项目适合参加人工智能或大数据赛道。

\textbf{6. 全国大学生数学建模竞赛}

每年9月举办，金融类题目（如贷款优化、投资组合、风险评估）是常见命题。适合有数学建模基础的团队参赛。

\subsection{参赛策略}

从选题到路演，竞赛准备通常分为五个阶段：

\textbf{第一阶段：选题（2周）}。选择有竞争力的选题是成功的第一步。竞赛选题应满足"三有"标准：有社会价值（解决真实痛点）、有技术亮点（有辨识度的创新）、有展示效果（能直观呈现）。建议从课程项目中提取最有竞争力的方向进行深化。

\textbf{第二阶段：组队（1周）}。竞赛团队需要"技术+业务+表达"三类人才。技术负责实现，业务负责行业洞察和需求分析，表达负责PPT和路演。最佳团队配置是2人技术+1人业务+1人表达。

\textbf{第三阶段：开发（3--4周）}。在课程项目基础上进行竞赛级升级：（1）增加数据量（真实数据优于模拟数据）；（2）完善功能（覆盖完整业务流程）；（3）优化体验（界面美观、操作流畅）；（4）增加亮点（可视化大屏、实时数据流等）。

\textbf{第四阶段：包装（1周）}。竞赛作品的"包装"极为重要：（1）商业计划书——从痛点分析、市场规模、解决方案、竞争优势、商业模式到团队介绍；（2）演示视频——3分钟，配解说，突出核心功能和创新点；（3）技术白皮书——详细的架构设计和技术决策说明。

\textbf{第五阶段：路演（准备+实战）}。路演是竞赛的决胜环节。准备要点：（1）反复演练，确保5分钟内完整展示；（2）准备3--5个评委可能追问的"坑"的应对方案；（3）团队配合演练，主答+补充配合默契；（4）路演时保持自信、简洁、有节奏感。

\subsection{从课程项目到竞赛作品}

课程项目和竞赛作品之间有一段"升级距离"，以下是关键的升级路径：

\begin{table}[H]
\centering
\small
\begin{tabular}{lll}
\toprule
\textbf{维度} & \textbf{课程项目} & \textbf{竞赛作品} \\
\midrule
数据 & 模拟数据 & 真实/半真实数据 \\
功能 & MVP可用 & 功能完整+亮点 \\
界面 & 基本可用 & 美观专业 \\
文档 & 实验报告 & 商业计划书+技术白皮书 \\
演示 & 课堂展示 & 路演视频+现场演示 \\
创新 & 至少一个点 & 多层次创新（技术+业务+方法） \\
受众 & 任课教师 & 评委+投资人+同行 \\
\bottomrule
\end{tabular}
\caption{课程项目到竞赛作品的升级}
\end{table}

\begin{casebox}[往届获奖项目分析]
\textbf{案例1：某高校"智能信贷风控平台"}（全国大学生金融科技创新大赛一等奖）

项目亮点：（1）真实数据——使用Kaggle公开的信贷违约数据集训练模型，准确率达87\%；（2）可解释性——每个拒贷决策附带SHAP值解释，满足监管要求；（3）可视化——风控大屏实时展示审批流程和风险热力图。

升级路径：从课程实验的"贷款违约预测"（纯模型训练）升级为完整的"风控平台"（包含审批流程、决策解释、可视化监控），增加了业务闭环和监管合规两个维度。

\textbf{案例2：某高校"银行ESG智能评估系统"}（"挑战杯"省级一等奖）

项目亮点：（1）多源数据——融合年报文本（NLP分析）、新闻舆情（情感分析）、碳排放数据（结构化分析）；（2）知识图谱——构建银行ESG本体，支持关联推理；（3）智能报告——一键生成符合监管要求的ESG评估报告。

升级路径：从课程实验的"ESG数据可视化"（纯展示）升级为"智能评估系统"（包含NLP分析、知识图谱、自动报告），增加了分析深度和自动化程度。

\textbf{案例3：某高校"普惠金融智能推荐引擎"}（"互联网+"省级银奖）

项目亮点：（1）用户画像——基于多维度数据构建农户/小微企业画像；（2）智能匹配——将用户需求与银行产品智能匹配；（3）社交传播——设计推荐奖励机制，实现低成本获客。

升级路径：从课程实验的"产品推荐"（基于规则匹配）升级为"智能推荐引擎"（包含用户画像、协同过滤、社交裂变），增加了推荐算法的深度和商业模式的可行性。
\end{casebox}

% ============================================================
\section{本章小结}
% ============================================================

本章围绕"如何做好课程综合项目"这一核心问题，从选题、管理、展示、转化四个维度进行了系统讲解：

\begin{enumerate}[leftmargin=2em]
  \item \textbf{选题}：遵循可行性、创新性、金融相关、技术适配四原则，12个可选方向覆盖不同难度
  \item \textbf{管理}：4周时间线、Feature Branch工作流、AI辅助的项目管理三件套
  \item \textbf{展示}：5分钟"问题→方案→演示→技术→总结"结构，功能完整性30\%+技术质量25\%的评分体系
  \item \textbf{转化}：从实验报告到学术论文的路径，三类论文选题方向和投稿目标
  \item \textbf{竞赛}：六大竞赛的参赛策略，从课程项目到竞赛作品的升级路径
\end{enumerate}

\begin{theorybox}[从"完成项目"到"创造价值"]
综合项目不仅是一次课程考核，更是一次将所学知识转化为实际价值的实践。好的项目应该能回答三个问题：解决了什么真实问题？用了什么创新方法？产生了什么实际效果？当你能清晰地回答这三个问题，项目就不再只是"作业"，而是一份值得展示的作品、一篇值得发表的论文、或一个值得参赛的创新方案。
\end{theorybox}

% ============================================================
\section{思考题}
% ============================================================

\begin{exercisebox}[本章思考题]
\begin{enumerate}[leftmargin=2em]
  \item 选择一个你感兴趣的项目方向，用选题自查清单逐项评估其可行性。
  \item 为你选择的项目方向设计4周时间线，标注关键里程碑和风险点。
  \item 如果你的团队只有2人，应如何调整角色分工和时间安排？
  \item 将课程中某个实验（如第11章智能客服）升级为竞赛作品，列出需要升级的3个维度及具体措施。
  \item 设计一个5分钟的答辩PPT大纲，包含每页的标题和核心内容。
  \item 查找一篇近两年发表的金融科技类学术论文，分析其创新点和研究方法，思考你的项目能否形成类似的学术贡献。
  \item 访问"互联网+"或"挑战杯"官网，了解最新的竞赛规则和评审标准，分析你的项目需要在哪些方面重点加强。
\end{enumerate}
\end{exercisebox}

% ============================================================
\section{AI辅助教学：OpenMAIC实践}
% ============================================================

\begin{theorybox}[清华大学OpenMAIC项目]
\textbf{OpenMAIC}（Open Multi-Agent Interactive Classroom，开放式多智能体交互课堂）是清华大学教育学院于济凡团队开发的开源AI教学平台（GitHub 18.6k Stars）。它采用多智能体编排技术，可将任何主题或文档一键转化为丰富的交互式课堂体验。
\end{theorybox}

\subsection{OpenMAIC核心功能}

\begin{table}[H]
\centering
\begin{tabular}{lp{8cm}}
\toprule
\textbf{功能} & \textbf{说明} \\
\midrule
演示幻灯片 & AI自动生成课件，支持语音讲解、聚光灯和激光笔效果 \\
随堂测验 & 自动生成单选/多选/简答题，AI实时评分反馈 \\
交互式模拟实验 & 可操作的HTML模拟场景（3D可视化、流程模拟、小游戏） \\
白板演示 & AI教师实时绘图讲解，支持公式推导、流程图绘制 \\
项目制学习(PBL) & 学生选择角色，与AI协作完成结构化项目 \\
\bottomrule
\end{tabular}
\caption{OpenMAIC核心功能}
\end{table}

\subsection{OpenMAIC在金融教学中的应用}

OpenMAIC特别适合金融科技课程的以下场景：

\textbf{场景一：概念讲解}。输入"商业银行数字化转型"，OpenMAIC自动生成包含AI教师讲解、AI同学讨论的互动课堂，学生可随时提问。

\textbf{场景二：案例分析}。输入"招商银行智慧零售实践"，系统生成交互式案例分析，学生可扮演不同角色（客户经理、风控官、产品经理）进行模拟决策。

\textbf{场景三：技术演示}。输入"MCP协议工作原理"，系统生成3D可视化演示，直观展示AI Agent与MCP Server的交互流程。

\textbf{场景四：项目答辩}。学生使用OpenMAIC为综合项目生成5分钟交互式课堂演示，作为答辩材料。

\subsection{快速上手}

\begin{lstlisting}[style=shell]
# 克隆仓库
git clone https://github.com/THU-MAIC/OpenMAIC.git
cd OpenMAIC

# 安装依赖
pnpm install

# 配置API Key（至少配置一个）
cp .env.example .env.local
# 编辑 .env.local，填入 OPENAI_API_KEY 或 ANTHROPIC_API_KEY

# 启动服务
pnpm dev
\end{lstlisting}

启动后访问 \texttt{http://localhost:3000}，输入学习主题即可开始。

\begin{notebox}[实践任务]
使用OpenMAIC为你的综合项目生成一个交互式课堂演示：
\begin{enumerate}[leftmargin=1.5em]
  \item 登录 \texttt{http://open.maic.chat}（托管版）或本地部署
  \item 输入你的项目主题（如"基于AI的银行智能客服系统"）
  \item 选择课堂模式：标准模式（课件+测验）或深度交互模式（模拟实验+编程）
  \item 导出为PPT或HTML，嵌入项目答辩材料
\end{enumerate}
\end{notebox}

\subsection{OpenMAIC与本课程工具链的关系}

\begin{table}[H]
\centering
\begin{tabular}{lll}
\toprule
\textbf{工具} & \textbf{定位} & \textbf{与OpenMAIC的关系} \\
\midrule
AI CLI（MiMo/Claude） & 代码生成与执行 & OpenMAIC可调用CLI生成交互式代码演示 \\
MCP协议 & 工具连接 & OpenMAIC的API可作为MCP Server接入 \\
Skill体系 & 领域知识 & OpenMAIC的Skill系统与本课程Skill兼容 \\
OpenMAIC & 教学呈现 & 将上述工具的成果转化为交互式课堂 \\
\bottomrule
\end{tabular}
\caption{OpenMAIC与课程工具链的协同}
\end{table}

\begin{notebox}[黄金组合]
CLI（执行层）+ Skill（知识层）+ MCP（连接层）+ OpenMAIC（呈现层）= 完整AI教学工具链。

使用CLI和Skill开发金融应用，通过MCP连接数据源，最后用OpenMAIC将成果转化为交互式课堂演示——这是本课程推荐的"从开发到展示"完整路径。
\end{notebox}
